% ─────────────────────────────────────────────────────────────────
%  Draft of §IV.C (Linear stability is necessary but not sufficient)
%  — slim version after retiring C1(b) (PDE-overlay panel) and C2
%  (T_PDE / T_LSA contrast).  Argument now leans on §IV.B's
%  spatially-generated slow manifold; C1(a) (analytic SS-existence
%  map with Whitney-cusp fold) and C3 (SNIC period-divergence scaling)
%  are the two empirical anchors.
% ─────────────────────────────────────────────────────────────────

\subsection{Linear stability is necessary but not sufficient}
\label{sec:lsa_vs_nl}

Section~\ref{sec:lcst_front_mechanism} showed that the LCST-front
attractor is a relaxation cycle whose slow manifold is generated by
the surface boundary condition together with the
$\varphi$-dependent diffusion barrier.  No spatially uniform reduction
can produce that manifold: at the surface the locus
$\mu(J,\theta)=\mu_b$ is enforced by the Robin BC, while in the
interior the same locus is irrelevant because the bulk equilibrium is
buffered by transport.  Linear stability of an available homogeneous
branch is therefore at best a \emph{necessary} condition for the PDE
to oscillate; it cannot be sufficient because the cycle traverses
two branches of a slow manifold that does not exist in the 0D
reduction.  This subsection makes the necessary-condition statement
quantitative on two complementary axes: (i)~where the homogeneous SS
even exists as a real root, and (ii)~how the cycle is destroyed at
the parameter-plane boundaries of the LCST-front regime.

\paragraph{Existence of the homogeneous swollen branch (analytic).}
The 0D steady-state system \eqref{eq:lsa_ss} admits an explicit
parameterization: at any pair $(J,\theta)$ with $J>\varphi_{p0}$,
the two SS conditions solve in closed form for the dependent
control parameters,
\begin{equation}
  S_\chi(J,\theta) =
    \frac{\mu_b - \mu_{\mathrm{base}}(J,\theta)}{\theta\,\varphi^2}, \qquad
  \Bi_T(J,\theta) = \frac{K(J,\theta)}{\theta\,(1+K/\Bi_c)},
  \label{eq:ss_param_inv}
\end{equation}
where $\mu_{\mathrm{base}}$ is the part of $\mu$ that excludes the
$S_\chi\theta\,\varphi^2$ coupling and
$K(J,\theta)=\Da\,J\,(1-\varphi)^{m_\mathrm{act}}\,T(\theta)$.  The
entire 0D SS surface is therefore the image of the $(J,\theta)$ plane
under \eqref{eq:ss_param_inv}, and the saddle--node fold locus in
$(\Bi_T, S_\chi)$ is exactly the image of the curve
$\det\partial(F_1,F_2)/\partial(J,\theta)=0$.  Numerically that locus
is a single connected component in $(J,\theta)$; under the map
\eqref{eq:ss_param_inv} it self-folds in $(\Bi_T,S_\chi)$ into two
branches meeting at a Whitney cusp at
$(\Bi_T,S_\chi)\!\simeq\!(0.43, 2.17)$ — a cusp catastrophe in the
homogeneous-SS surface.

Counting the real roots of \eqref{eq:lsa_ss} cell-by-cell on a
$360\times 240$ grid in $(\Bi_T, S_\chi)$ via the closed-form
inverse \eqref{eq:ss_param_inv} (one sign-change scan per cell along
$J$, exact and ${\sim}10^{2}\times$ faster than fsolve) gives
Fig.~\ref{fig:fig_homog_ss}.  The fraction of the displayed plane
occupied by each multiplicity is summarised below
\begin{center}
\begin{tabular}{lr}
\hline
Region & Area fraction \\
\hline
swollen branch only & $37\%$ \\
collapsed branch only & $51\%$ \\
bistable (3 roots: 2 stable + 1 saddle) & $11\%$ \\
no root (off the displayed window) & $0.3\%$ \\
\hline
\end{tabular}
\end{center}
exactly $0\%$ of cells have $2$ roots: this odd-only multiplicity is
forced by the fold topology (any horizontal line in
$(\Bi_T, S_\chi)$ at fixed $S_\chi$ crosses the two fold arcs an even
number of times, giving root counts $1, 3, 5, \ldots$).  The default
working point $(\Bi_T=0.10,S_\chi=1.0)$ and the canonical C-cell
$(\Bi_T=0.059, S_\chi=1.80)$, both used elsewhere in this section,
sit deep inside the ``collapsed-only'' region: at neither point does
the swollen homogeneous SS that one would linearise about exist as a
real root, so a 0D LSA about the swollen branch is structurally
unavailable.  The PDE nevertheless finds the LCST-front cycle there.
The mismatch is the
necessary-but-not-sufficient statement made geometric: the
spatially-generated slow manifold of
Sec.~\ref{sec:lcst_front_mechanism} continues to exist past the
saddle--node fold of the homogeneous problem because its swollen
branch is enforced \emph{at the surface} by the Robin BC, not
required to coexist with bulk reactant balance.

\paragraph{SNIC scaling at the upper-$S_\chi$ exit.}
At the upper-$S_\chi$ boundary of the LCST-front regime
[Fig.~\ref{fig:fig4_phase}(e)], the cycle period diverges as
$S_\chi\to S_\chi^{(c)}$ from below.  We test the period scaling
along a 1D slice through this boundary against the two canonical
infinite-period scenarios
\begin{equation}
  T_\mathrm{SNIC}(S_\chi)\propto (S_\chi^{(c)}-S_\chi)^{-1/2},
  \qquad
  T_\mathrm{hom}(S_\chi)\propto -\ln(S_\chi^{(c)}-S_\chi),
  \label{eq:snic_vs_hom}
\end{equation}
using the same 1D scan
[Fig.~\ref{fig:fig_snic_scaling}].  A two-parameter least-squares fit
of $T_\mathrm{PDE}(S_\chi)$ on the bracket
$S_\chi\in[1.10, 1.32]$ yields a SNIC residual sum of squares
$\mathrm{RSS}_\mathrm{SNIC}=0.068$ versus
$\mathrm{RSS}_\mathrm{hom}=0.152$, with the SNIC fit giving the
critical value $S_\chi^{(c)}\!\approx\!1.34$.  The SNIC inverse-square-
root law captures the data more than twice as well as the homoclinic
log law, consistent with the cycle being destroyed by the collision
of a saddle and a node on the limit cycle rather than by a homoclinic
orbit hitting a saddle off the cycle.  In either reading the cycle's
disappearance is a global topological event in the surface
$(\theta_\mathrm{surf}, J_\mathrm{surf})$ slow-manifold plane and not
a local linear bifurcation of any single homogeneous branch — the
expected signature of a relaxation oscillator on a spatially
generated slow manifold.

\paragraph{Summary.}
The 0D linear stability analysis carries a precise but limited
informational content: it identifies the local stability of available
homogeneous branches, but the swollen homogeneous branch fails to
exist over $\sim\!51\%$ of the parameter plane covered here, and even
where it does the cycle that the PDE realises is generically not the
small-amplitude continuation of any homogeneous Hopf bifurcation.
The dominant nonlinear behaviour of this thermo-LCST gel is a
relaxation oscillator on a spatially generated slow manifold, born
and destroyed by global topological events
(Sec.~\ref{sec:lcst_front_mechanism} and Eq.~\ref{eq:snic_vs_hom}).
Whether the same parameter plane harbours a hidden second attractor
not visible to either linear analysis or default-IC numerical scans
is settled in Sec.~\ref{sec:multistability}.
